\documentclass{article}
\usepackage{amsmath, amssymb, amsthm}
\usepackage{hyperref}
\usepackage{geometry}
\geometry{margin=1in}

\title{Erd\H{o}s Problem \#679}
\date{Last updated: 2025-08-31}
\author{Source: erdosproblems.com}

\begin{document}
\maketitle

\noindent\textbf{Status:} OPEN \quad \textbf{No prize}

\noindent\textbf{Tags:} number theory

\medskip
\noindent\textbf{Problem Statement:}

Let $\epsilon>0$ and $\omega(n)$ count the number of distinct prime factors of $n$. Are there infinitely many values of $n$ such that\[\omega(n-k) < (1+\epsilon)\frac{\log k}{\log\log k}\]for all $k<n$ which are sufficiently large depending on $\epsilon$ only?

Can one show the stronger version with\[\omega(n-k) < \frac{\log k}{\log\log k}+O(1)\]is false?

\bigskip
\noindent\textbf{Remarks:}

One can ask similar questions for $\Omega$, the number of prime factors with multiplicity, where $\log k/\log\log k$ is replaced by $\log_2k$.

In the comments DottedCalculator has disproved the second stronger version, proving that in fact for all large $n$ there exists $k<n$ such that\[\omega(n-k)\geq \frac{\log k}{\log\log k} + c\frac{\log k}{(\log\log k)^2}\]for some constant $c>0$.

Lau \cite{La26} has proved that there exists a constant $C>0$ such that, for infinitely many $n$,\[\omega(n-k) \leq \Omega(n-k)\leq C\log k\]for all $1<k<n$. Lau conjectures that, for $\omega$, this can only be improved in the constant factor, and thus the original question has a negative answer.

See also [248], [413], and [1203].

\bigskip
\noindent\textbf{References:}

[La26] C. F. Lau, On the number of prime factors of consecutive integers. arXiv:2604.15042 (2026).

\bigskip
\noindent\small{Source: \url{https://www.erdosproblems.com/679}}
\end{document}
